\section{中都、河南与边区的不同负担}

1153年迁都中都、1156年集结南侵军资、1161年渡江受挫，留下的是不同区域被卷入财政与军事的方式。中都所需的营建、仓储与官署不能由一处城址推算全国劳役；河南和边区的驻军、转运与市场又有不同的道路和供给条件。《金史》制度叙述适合建立政策线索，碑志、窑冶遗存、钱币与城址才可能检验地方生产和迁徙的部分痕迹。

因此，不应把金的地方社会写为“女真”与“汉人”两种统一经验。官号、籍属、语言、职业、家族与居地记录的事实并不相同。金宋边界的和议和战事，既改变军队和税源，也影响商人、工匠与被迁徙者；但现有材料密度不均，尤其不能用中都或一处边城取代北方各路。
